\subsection{The three-cycle variation}
\label{rw:lower:preview}

The reversible center also gives a direct route to the lower bound.
Keep the original coefficients $B=0$, $m=\epsilon>0$ and start at
$a=c=d=0$. Its bounded period annulus has the upper saddle and the
regular lower paths prepared in Section~\ref{rw:transitions}.
Use the actual upper ratio $r=1+\kappa$ as a coefficient coordinate.
Section~\ref{rw:lower-bound-section} constructs a single coefficient
curve
\[
 (\bar a,\kappa,\nu)
       =(h_\epsilon\vartheta,\vartheta,\beta_\epsilon\vartheta),
 \qquad h_\epsilon\longrightarrow\tfrac32,
\]
with exponentially small $\beta_\epsilon$ and positive
$\delta_\epsilon,S_\epsilon$. The ratio variation produces the
logarithmic term in the closing difference, and the line-breaking
coefficient fixes its constant term. For the actual central-height
closing function the resulting profile is
\[
 \lim_{\epsilon\downarrow0}\lim_{\vartheta\downarrow0}
 \frac{F(\delta_\epsilon x;\vartheta)}{\vartheta S_\epsilon}
 =P(x):=\frac{x^2}{2}-x\log x-x-\frac1{20},
\]
with $C^1$ convergence on compact subintervals of $(0,10]$.
The function $P$ has signs $-,+,-,+$ at $10^{-3},10^{-1},1,10$;
its derivative $x-\log x-2$ has exactly two simple zeros.
It therefore has three simple positive roots.

For each fixed sufficiently small $\epsilon$, one sufficiently small
positive $\vartheta$ preserves all three roots on three disjoint
compact bands of real center ovals. These are three cycles of the
same original field. Proposition~\ref{rw:lower-variation} proves
the first-variation profile; Theorem~\ref{rw:three-simple-cycles}
proves the fixed-scale perturbation limit, simultaneous persistence,
positive complete returns and convergence of each cycle to the entire
graphic.
The forward return also determines their alternating stability in
Corollary~\ref{rw:lower:stability}. This construction uses the
reversible center and its prepared paths independently of the upper
count.
